\begin{proof}[Proof of Proposition~\ref{prop:planning-solution}]
    As stated in the text, the present value~\eqref{eq:planning-present-value} depends on~$n$ only via the residual variance~$R_n$ (see Footnote~\ref{fn:planning-present-value}).
    Let~$\Psi(R_n)$ denote this present value when the agent chooses the precisions~$p_n,p_{n+1},\ldots$ that solve~\eqref{eq:planning-problem}.
    Then
    %
    \begin{align}
        \Psi(R_n)
        &= \min_{p_n\ge0}\left\{\left(\frac{1}{R_n}+p_n\right)^{-1}+cp_n+\delta\Psi\left(\left(\frac{1}{R_n}+p_n\right)^{-1}+\sigma^2\right)\right\} \notag \\
        &= \min_{V_n\in(0,R_n]}\left\{V_n+c\left(\frac{1}{V_n}-\frac{1}{R_n}\right)+\delta\Psi\left(V_n+\sigma^2\right)\right\}, \label{eq:planning-solution-bellman}
    \end{align}
    %
    where the posterior variance
    %
    \begin{align*}
        V_n
        &\equiv \Var(\theta(t_n)\mid\Hcal_n) \\
        &= \left(\frac{1}{R_n}+p_n\right)^{-1}
    \end{align*}
    %
    is an invertible function of~$p_n$ given~$R_n$ and so the change of variables is well-defined.
    We want to find a function~$\Psi:(0,\infty)\to(0,\infty)$ satisfying the Bellman equation~\eqref{eq:planning-solution-bellman}.
    This is equivalent to finding a function~$\Phi:(0,\infty)\to(0,\infty)$ mapping residual variances to optimal posterior variances.
    
    Assume~$\Psi$ and~$\Phi$ are continuous and piecewise differentiable.
    (One can verify these properties using arguments provided in Section~3.3 of \citet{Stokey-Lucas-1989-}.)

    Now~$V_n\le R_n$ with equality if and only if~$p_n=0$.
    But choosing~$p_n=0$ is optimal if and only if the cost of buying information exceeds its value, which occurs precisely when~$R_n$ is sufficiently small.
    So there exists~$R>0$ such that~$\Phi(R_n)=R_n$ for all~$R_n<R$ and~$\Phi(R_n)<R_n$ for all~$R_n>R$.

    Suppose~$R_n>R$.
    Then the constraint~$V_n\le R_n$ on the RHS of~\eqref{eq:planning-solution-bellman} is non-binding, and so
    %
    \begin{align*}
        \Psi'(R_n)
        &= \parfrac{}{R_n}\min_{V_n\in(0,R_n]}\left\{V_n+c\left(\frac{1}{V_n}-\frac{1}{R_n}\right)+\delta\Psi\left(V_n+\sigma^2\right)\right\} \\
        &= \frac{c}{R_n^2}
    \end{align*}
    %
    by the envelope theorem.
    So the optimal posterior variance satisfies the first-order condition
    %
    \begin{align*}
        0
        &= \parfrac{}{V_n}\left\{V_n+c\left(\frac{1}{V_n}-\frac{1}{R_n}\right)+\delta\Psi\left(V_n+\sigma^2\right)\right\} \\
        &= 1-\frac{c}{V_n^2}+\delta\Psi'(V_n+\sigma^2) \\
        &= 1-c\left(\frac{1}{V_n^2}-\frac{\delta}{(V_n+\sigma^2)^2}\right),
    \end{align*}
    %
    which is uniquely solved by~$V$.
    Thus~$\Phi(R_n)=V$ for all~$R_n>R$, from which it follows that~$R=V$ since~$\Phi$ is continuous.
    So the solution to~\eqref{eq:planning-problem} induces residual variances~$R_n,R_{n+1},\ldots$ that satisfy
    %
    \begin{align*}
        R_{i+1}
        &= \Phi(R_i)+\sigma^2 \\
        &= \min\{R_i+\sigma^2,V+\sigma^2\}
    \end{align*}
    %
    for each~$i\ge n$.
    Therefore, if the agent chooses the precisions that solve~\eqref{eq:planning-problem} for each~$n\in\N$, then for each~$n>1$ we have
    %
    \begin{align*}
        R_n
        &= \min\{R_{n-1}+\sigma^2,V+\sigma^2\} \\ 
        % &= \min\{\min\{R_{n-2}+\sigma^2,V+\sigma^2\}+\sigma^2,V+\sigma^2\} \\
        % &= \min\{R_{n-2}+2\sigma^2,V+2\sigma^2,V+\sigma^2\} \\
        % &= \min\{R_{n-3}+3\sigma^2,V+3\sigma^2,V+2\sigma^2,V+\sigma^2\} \\
        &\hspace{0.6em} \vdots \\
        &= \min\{R_1+\sigma^2(n-1)\}\cup\left\{V+\sigma^2i\right\}_{i=1}^{n-1} \\
        &= \min\{\sigma_0^2+\sigma^2t_n,V+\sigma^2\}.
    \end{align*}
    %
    It follows that, for each~$n\in\N$, the optimal precision
    %
    \begin{align*}
        \frac{1}{\min\{R_n,V\}}-\frac{1}{R_n}
        &= \frac{1}{\min\{\sigma_0^2+\sigma^2t_n,V\}}-\frac{1}{\min\{\sigma_0^2+\sigma^2t_n,V+\sigma^2\}} \\
        &= \begin{cases}
            0 & \text{if}\ \sigma_0^2+\sigma^2t_n<V \\
            \frac{1}{V}-\frac{1}{\sigma_0^2+\sigma^2t_n} & \text{if}\ V\le\sigma_0^2+\sigma^2t_n<V+\sigma^2 \\
            \frac{1}{V}-\frac{1}{V+\sigma^2} & \text{if}\ V+\sigma^2\le\sigma_0^2+\sigma^2t_{n-1}
        \end{cases} \\
        &= p_n^\dag
    \end{align*}
    %
    induces posterior variance
    %
    \begin{align*}
        \Var(\theta(t_n)\mid\Hcal_n)
        &= \min\{R_n,V\} \\
        &= \min\{\sigma_0^2+\sigma^2t_n,V\}.\qedhere
    \end{align*}
\end{proof}
